% ============================================================
% Master Thesis Skeleton (Balanced Theory + Practice)
% Topic: Lightweight semantics for job posts & CSV uploads
% ============================================================

\documentclass[12pt,a4paper,oneside]{report}

% ----------------------------
% Packages (keep it reasonable)
% ----------------------------
\usepackage[a4paper,margin=1in]{geometry}
\usepackage{setspace}          % line spacing
\usepackage{graphicx}          % figures
\usepackage{booktabs}          % nice tables
\usepackage{longtable}         % long tables across pages
\usepackage{amsmath,amssymb}   % math
\usepackage{hyperref}          % links and references
\usepackage{url}               % url formatting
\usepackage{caption}           % captions
\usepackage{subcaption}        % subfigures
\usepackage{float}             % figure placement
\usepackage{enumitem}          % nicer lists
\usepackage{csquotes}          % recommended for biblatex
\usepackage{listings}          % code listings
\usepackage{xcolor}            % code colors
\usepackage{algorithm}
\usepackage{algpseudocode}     % algorithms pseudocode
\usepackage{glossaries}        % glossary / acronyms
\usepackage{biblatex}          % bibliography management

% ----------------------------
% Bibliography
% ----------------------------
\addbibresource{references.bib}

% ----------------------------
% Formatting
% ----------------------------
\onehalfspacing
\hypersetup{
  colorlinks=true,
  linkcolor=blue,
  citecolor=blue,
  urlcolor=blue,
  pdftitle={Master Thesis Title},
  pdfauthor={Your Name}
}

% ----------------------------
% Listings configuration
% ----------------------------
\lstset{
  basicstyle=\ttfamily\small,
  breaklines=true,
  frame=single,
  numbers=left,
  numberstyle=\tiny,
  keywordstyle=\color{blue},
  commentstyle=\color{gray},
  stringstyle=\color{teal}
}

% ----------------------------
% Acronyms / Glossary (optional)
% ----------------------------
\makeglossaries
% Example:
% \newacronym{nlp}{NLP}{Natural Language Processing}
% \newacronym{ir}{IR}{Information Retrieval}
% \newacronym{ner}{NER}{Named Entity Recognition}
% \newacronym{bm25}{BM25}{Best Matching 25}

% ----------------------------
% Thesis metadata (edit these)
% ----------------------------
\newcommand{\ThesisTitle}{Lightweight Semantic Structuring of Internship Listings from Web Scraping and CSV Uploads}
\newcommand{\AuthorName}{Mouad bensaid}
\newcommand{\DegreeName}{Master of Science (or your program)}
\newcommand{\UniversityName}{Your University}
\newcommand{\DepartmentName}{Your Department}
\newcommand{\SupervisorName}{Supervisor Name}
\newcommand{\SubmissionDate}{February 2026}

% ----------------------------
% Title page
% ----------------------------
\begin{document}

\begin{titlepage}
  \centering
  {\Large \UniversityName \par}
  \vspace{0.5cm}
  {\large \DepartmentName \par}
  \vspace{2.0cm}
  {\huge \bfseries \ThesisTitle \par}
  \vspace{2.0cm}
  {\Large \AuthorName \par}
  \vspace{0.5cm}
  {\large \DegreeName \par}
  \vspace{1.5cm}
  {\large Supervisor: \SupervisorName \par}
  \vfill
  {\large \SubmissionDate \par}
\end{titlepage}

% ----------------------------
% Front matter
% ----------------------------
\pagenumbering{roman}

\chapter*{Abstract}
% 150--300 words. Problem, method, key results, contributions.
\addcontentsline{toc}{chapter}{Abstract}

\chapter*{Résumé / ملخص}
% Optional depending on university requirements.
\addcontentsline{toc}{chapter}{Résumé / ملخص}

\chapter*{Acknowledgements}
\addcontentsline{toc}{chapter}{Acknowledgements}


\addcontentsline{toc}{chapter}{Declaration}

\tableofcontents
\listoffigures
\listoftables

\printglossaries

% ============================================================
% Main matter
% ============================================================
\clearpage
\pagenumbering{arabic}

% ============================================================
% Part I: Theoretical Background (Aim ~50%)
% ============================================================
\part{Theoretical Background and Related Work}

\chapter{Introduction}
\section{Context and Motivation}
% Internship/job listings are noisy, heterogeneous, and often semi-structured.
% Motivation: better search, filtering, matching, analytics.

\section{Problem Statement}
% Define: input sources (scraped posts + user CSV uploads), target: semantic representation.
% Define constraints: low compute, low latency, low cost; limited/optional LLM.

\section{Research Questions}
\begin{itemize}[noitemsep]
  \item RQ1: How can scraped job postings and CSV uploads be mapped into a structured semantic schema using lightweight methods?
  \item RQ2: Which combination of rule-based extraction and small-model classification yields the best cost/quality trade-off?
  \item RQ3: How does semantic enrichment improve retrieval quality compared to keyword-only baselines?
\end{itemize}

\section{Contributions}
\begin{itemize}[noitemsep]
  \item A unified schema for internships that supports filtering, ranking, and analytics.
  \item A lightweight enrichment pipeline (rules + dictionaries + small ML).
  \item An evaluation framework and ablation study comparing components.
\end{itemize}

\section{Thesis Structure}
% Briefly describe each part and chapter.

\chapter{Background: Information Extraction and Semantics}
\section{Semi-Structured Web Data and Job Postings}
% Discuss typical noise, layout variability, duplicates, expiry.

\section{Information Extraction (IE) Fundamentals}
% Named entities, relations, slot filling, pattern-based extraction.

\section{Semantic Modeling}
\subsection{Schemas vs Ontologies vs Taxonomies}
% Define each and why a schema-first approach fits the problem.

\subsection{Skill and Occupation Taxonomies}
% Introduce ESCO / O*NET and explain how you use them (as vocabulary, not full ontology).

\section{Normalization and Canonicalization}
% Titles normalization, synonyms, deduplication, standard fields.

\chapter{Background: Text Representation and Retrieval}
\section{Keyword Retrieval}
% BM25 and inverted indexes; strengths/weaknesses for noisy job text.

\section{Vector Representations}
% TF-IDF, word embeddings, sentence embeddings; lightweight options.

\section{Hybrid Retrieval}
% Combining lexical + semantic vectors; practical justification.

\section{Efficiency Considerations}
% Latency, memory, inference cost; caching; indexing strategies.

\chapter{Related Work}
\section{Job Posting Analytics and Skill Extraction}
% Prior approaches: dictionaries, weak supervision, ML, LLM-based.

\section{Ontology/Taxonomy Alignment in HR/Recruitment}
% Mapping job text to standardized labels.

\section{LLM-based Pipelines vs Lightweight Pipelines}
% Discuss trade-offs: cost, explainability, error modes.

% ============================================================
% Part II: Practical Design and Implementation (Aim ~50%)
% ============================================================
\part{System Design, Implementation, and Evaluation}

\chapter{Methodology and System Overview}
\section{Design Goals and Constraints}
\begin{itemize}[noitemsep]
  \item Low computational cost per document
  \item Robustness to noisy text and HTML variability
  \item Explainability: traceable extraction decisions
  \item Extensibility: adding skills/labels without retraining everything
\end{itemize}

\section{Proposed Architecture}
% Include a system diagram figure.
\begin{figure}[H]
  \centering
  % \includegraphics[width=0.9\textwidth]{figures/architecture.pdf}
  \caption{High-level architecture of the data ingestion and semantic enrichment pipeline.}
  \label{fig:architecture}
\end{figure}

\section{Data Flow}
% Scraper -> Cleaner -> Extractor -> Normalizer -> Store -> Search API

\chapter{Data Collection and Preprocessing}
\section{Web Scraping Strategy}
% Sources, crawling frequency, compliance, deduplication.

\section{HTML Parsing and Text Extraction}
% Boilerplate removal, section detection (Responsibilities/Requirements).

\section{CSV Upload Handling}
\subsection{Column Mapping Heuristics}
% Header similarity + content-based heuristics.

\subsection{Validation and Sanitization}
% Required columns, encoding issues, missing values.

\section{Dataset Description}
% Volume, time span, language(s), distribution of fields.

\chapter{Semantic Schema and Label Space}
\section{Target Schema}
% Provide JSON schema and rationale.
\begin{lstlisting}[language=json,caption={Example target schema (simplified).},label={lst:schema}]
{
  "job_id": "...",
  "title_raw": "...",
  "title_normalized": "...",
  "seniority": "intern",
  "employment_type": "internship",
  "location": {"city": "...", "region": "...", "country": "CA", "mode": "hybrid"},
  "skills": [{"name": "python", "weight": 0.9}, {"name": "sql", "weight": 0.7}],
  "compensation": {"min": null, "max": null, "currency": "CAD", "period": "hour"},
  "source": {"url": "...", "site": "..."},
  "timestamps": {"scraped_at": "...", "posted_at": "..."}
}
\end{lstlisting}

\section{Canonical Labels}
% Define canonical titles/domains and how you curate them.

\section{Synonyms and Alias Tables}
% Example: js -> javascript, node -> node.js

\chapter{Lightweight Semantic Enrichment Pipeline}
\section{Layer A: Rule-Based Extraction}
\subsection{Patterns for Compensation, Dates, Location}
% Regex patterns and validation.

\subsection{Seniority and Employment Type}
% Keywords, disambiguation rules.

\subsection{Skill Matching with Lexicon}
% Dictionary-driven phrase matching; handling multi-word skills.

\begin{algorithm}[H]
\caption{Lexicon-based skill extraction (overview)}
\label{alg:skill_extraction}
\begin{algorithmic}[1]
\Require Text $T$, skill lexicon $L$ with synonyms and canonical forms
\Ensure Extracted skills $S$
\State $T' \gets$ normalize($T$)  \Comment{lowercase, punctuation, unicode, etc.}
\State $M \gets$ find\_matches($T', L$) \Comment{phrase matching}
\State $S \gets$ canonicalize\_and\_dedupe($M$)
\State $S \gets$ score\_skills($S, T'$) \Comment{weights via heuristics}
\State \Return $S$
\end{algorithmic}
\end{algorithm}

\section{Layer B: Small-Model Mapping}
\subsection{Title Normalization and Classification}
% TF-IDF + linear classifier OR small embedding + nearest-label mapping.

\subsection{Domain Classification (Optional)}
% Describe features and training procedure.

\subsection{Confidence Scoring}
% How you decide whether output is reliable.

\section{Layer C: Optional Fallback (Budgeted)}
% If you include an LLM at all: only for low-confidence cases.
% Must include a strict gating strategy and cost accounting.

\chapter{Storage, Indexing, and Search}
\section{Data Storage Model}
% DB tables/collections; versioning; auditability.

\section{Search Index Design}
% Fields: normalized, faceted filters; plus optional vector index.

\section{Hybrid Retrieval Strategy}
% BM25 + vectors, reranking rules.

\section{API Endpoints}
% /search, /jobs/{id}, /upload/csv, /stats

\chapter{Evaluation Setup}
\section{Evaluation Objectives}
% Extraction accuracy, classification accuracy, retrieval quality, cost/latency.

\section{Ground Truth Construction}
% Annotation guidelines; inter-annotator agreement (if possible).
% If limited time: small labeled set + spot checks.

\section{Metrics}
\subsection{Field-Level Extraction}
% Precision/Recall/F1 for skills, locations, seniority, etc.

\subsection{Classification}
% Accuracy, macro-F1; confusion matrix.

\subsection{Retrieval}
% nDCG@k, Precision@k, Recall@k, MRR.

\subsection{Efficiency}
% Latency per doc, throughput, memory, $ cost estimate.

\section{Baselines and Ablations}
\begin{itemize}[noitemsep]
  \item Baseline 1: keyword-only search (no enrichment)
  \item Baseline 2: rules-only extraction
  \item Baseline 3: rules + small-model mapping
  \item Ablations: remove title normalization / remove synonym table / remove section weighting
\end{itemize}

\chapter{Results}
\section{Extraction Results}
% Tables + error breakdown.

\section{Classification Results}
% Confusion matrix figure + discussion.

\section{Retrieval Results}
% Compare hybrid vs keyword-only.

\section{Efficiency and Cost}
% Latency histograms, per-document compute, scaling discussion.

\chapter{Discussion}
\section{Error Analysis}
% Common failure cases: ambiguous titles, missing info, skill synonyms, noisy HTML.

\section{Trade-offs: Accuracy vs Cost}
% Where lightweight wins/loses compared to heavy LLM.

\section{Threats to Validity}
% Sampling bias, annotation bias, generalizability, source drift.

\chapter{Ethical, Legal, and Privacy Considerations}
\section{Scraping Compliance and Respect for Terms}
% Robots.txt, rate limiting, attribution, opt-out mechanisms.

\section{Privacy and User Uploads}
% CSV handling, PII minimization, retention policy.

\section{Fairness and Bias}
% Biased postings, biased labels, mitigation strategies.

\chapter{Conclusion and Future Work}
\section{Summary of Findings}
% Answer RQs succinctly.

\section{Future Directions}
\begin{itemize}[noitemsep]
  \item Multilingual extraction
  \item Better section segmentation
  \item Self-training / distillation from limited LLM labeling
  \item Continuous monitoring for taxonomy drift
\end{itemize}

% ============================================================
% Back matter
% ============================================================
\clearpage
\printbibliography

% ============================================================
% Appendices
% ============================================================
\appendix

\chapter{Annotation Guidelines (if you label data)}
% Definitions, examples, edge cases.

\chapter{Regex and Heuristic Rules}
% Document patterns; version them.

\chapter{Schema Documentation}
% Field definitions, allowed values, examples.

\chapter{Skill Lexicon and Synonyms (Excerpt)}
% Provide a representative subset; full list can be linked or attached.

\end{document}
